%% ============================================================
%% AULA 1 — Introdução: redes neurais como aproximadores universais
%% GBC073 — Inteligência Computacional (FACOM/UFU)
%%
%% Versão sucinta: roteiro do apresentador (poucas palavras,
%% mesmas definições e figuras do aula01.tex)
%%
%% Espinha dorsal: CMU 11-785 (S26), Lecture 1 — traduzida e
%% adaptada: conexionismo, perceptron, portas lógicas, geometria
%% da separabilidade, XOR, universalidade e o papel da
%% profundidade. PyTorch desde o primeiro dia.
%% Complementada com material do MIT 6.S191 (Amini, ed. 2026).
%%
%% Compilar:  latexmk -xelatex aula01-sucinta.tex   (duas passadas!)
%% ============================================================
\documentclass[aspectratio=169, 11pt]{beamer}

\makeatletter
\def\input@path{{./}{../}}
\makeatother

\usetheme{InteligenciaComputacional}   % ANTES do babel: fontspec vem primeiro
\usepackage{babel}
\babelprovide[import, main]{brazilian}
\usepackage{booktabs}
\usepackage{listings}
\lstset{basicstyle=\ttfamily\footnotesize, showstringspaces=false}

\title[Aula 1 — Introdução]{Aula 1: Introdução — redes neurais como aproximadores universais}
\subtitle[GBC073 — Inteligência Computacional]{Conexionismo \textbullet\ Perceptron \textbullet\ Universalidade \textbullet\ O papel da profundidade}
\author[Prof.\ Marcelo Albertini]{Prof.\ Marcelo Keese Albertini}
\institute{GBC073 — Inteligência Computacional \textbullet\ Universidade Federal de Uberlândia}
\date{2026}

% foto de fonte aberta com fallback:
% \fotoslot{arquivo}{largura-max}{altura-max}{legenda}{crédito}
% se imagens/arquivo não existir, desenha um placeholder e o PDF compila igual.
\newcommand{\fotoslot}[5]{%
  \begin{center}%
  \IfFileExists{imagens/#1}{%
    \includegraphics[width=#2, height=#3, keepaspectratio]{imagens/#1}\par
  }{%
    \begin{tikzpicture}
      \draw[icCinza!55, dashed, thick, rounded corners=3pt]
        (0,0) rectangle ({#2},{#3});
      \node[align=center, text=icCinza, font=\scriptsize]
        at ({0.5*#2},{0.5*#3})
        {foto: \texttt{imagens/#1}\\[1pt] rode \texttt{imagens/baixar-imagens.sh}};
    \end{tikzpicture}\par
  }%
  {\scriptsize #4\par}%
  {\tiny\color{icCinza} #5\par}%
  \end{center}}

% mini-gráfico de separabilidade: pontos do plano booleano
% #1 = rótulo do título; corpo é passado como pontos já desenhados
\newcommand{\eixosbool}{%
  \draw[-{Stealth}, icCinza!80] (-0.45,0) -- (1.9,0) node[below, font=\scriptsize] {$x_1$};
  \draw[-{Stealth}, icCinza!80] (0,-0.45) -- (0,1.9) node[left, font=\scriptsize] {$x_2$};
  \foreach \p in {0,1.3}{\draw[icCinza!60] (\p,-0.05) -- (\p,0.05); \draw[icCinza!60] (-0.05,\p) -- (0.05,\p);}
  \node[font=\tiny, text=icCinza] at (1.3,-0.25) {$1$};
  \node[font=\tiny, text=icCinza] at (-0.25,1.3) {$1$};
}
\newcommand{\pzero}[2]{\node[circle, draw=icCiano, thick, fill=white, inner sep=2.2pt] at (#1,#2) {};}
\newcommand{\pum}[2]{\node[circle, draw=icMagenta, fill=icMagenta, inner sep=2.2pt] at (#1,#2) {};}

\begin{document}

% ------------------------------------------------------------ capa
\begin{frame}[plain, noframenumbering]
  \titlepage
\end{frame}

% ============================================================
\section{Por que esta disciplina, agora}

% ------------------------------------------------------------
\begin{frame}{O que ``inteligência computacional'' significa em 2026}
  \framesubtitle{A ementa lista três paradigmas; a história já votou entre eles}

  \begin{itemize}\setlength{\itemsep}{0.3ex}
    \item<1-> {\color{icNeural}\textbf{Redes neurais artificiais}} — de perceptrons
          a \emph{Transformers}: a \alert{tecnologia dominante} em visão,
          linguagem, fala, biologia e jogos
    \item<2-> {\color{icEvol}\textbf{Computação evolutiva}} — otimização
          \emph{sem gradiente}: onde ainda é competitiva (e por que
          perdeu o centro do palco)
    \item<3-> {\color{icFuzzy}\textbf{Sistemas nebulosos}} — lógica com graus de
          verdade: semente histórica das \emph{relaxações diferenciáveis}
  \end{itemize}

  \onslide<4->{%
  \begin{ideiabox}
    Fio condutor do semestre: \textbf{tudo que é diferenciável se treina por
    gradiente}. Os outros paradigmas viram contraponto: o que fazer quando o
    gradiente \emph{não} existe.
  \end{ideiabox}}
\end{frame}

% ------------------------------------------------------------
\begin{frame}{Três círculos concêntricos}
  \framesubtitle{Onde o deep learning mora — e o que cada fronteira acrescenta}

  \begin{columns}[T]
    \begin{column}{0.4\textwidth}
      \begin{center}
        \scalebox{0.8}{%
        \begin{tikzpicture}
          \fill[icNeural!12]    (0,0) circle (1.55);
          \fill[icAzulClaro!18] (0,0) circle (1.05);
          \fill[icMagenta!22]   (0,0) circle (0.55);
          \draw[icNeural, very thick]    (0,0) circle (1.55);
          \draw[icAzulClaro, very thick] (0,0) circle (1.05);
          \draw[icMagenta, very thick]   (0,0) circle (0.55);
          \node[font=\small, text=icNeural] at (0,1.28) {IA};
          \node[font=\small, text=icAzulClaro] at (0,0.8) {AM};
          \node[font=\small\bfseries, text=icMagenta] at (0,0) {DL};
        \end{tikzpicture}}
      \end{center}
    \end{column}
    \begin{column}{0.56\textwidth}
      \begin{itemize}\setlength{\itemsep}{0.35ex}
        \item {\color{icNeural}\textbf{Inteligência artificial}} — técnicas que
              imitam comportamento humano
        \item {\color{icAzulClaro}\textbf{Aprendizado de máquina}} — aprender
              sem ser explicitamente programado
        \item {\color{icMagenta}\textbf{Aprendizado profundo}} — padrões dos
              dados brutos, redes de múltiplas camadas
      \end{itemize}
    \end{column}
  \end{columns}

  \begin{ideiabox}
    Frase-resumo: ensinar o computador a aprender a tarefa a partir de dados
    brutos — em vez de receber a regra pronta.
  \end{ideiabox}
\end{frame}

% ------------------------------------------------------------
\begin{frame}{O momento em que estamos}
  \framesubtitle{Este não é um curso de história — mas a história acabou de nos dar razão}

  \begin{columns}[T]
    \begin{column}{0.60\textwidth}
      \begin{itemize}\setlength{\itemsep}{0.35ex}
        \item \textbf{2024:} Nobel de \emph{Física} para Hopfield e Hinton
              (redes neurais); de \emph{Química} para o AlphaFold
        \item \textbf{Modelos de linguagem} escrevem código, provam teoremas
              e passam em exames profissionais
        \item \textbf{Difusão} gera imagens, áudio, vídeo e moléculas
        \item Tudo é, no fundo, \alert{uma única receita}: redes profundas +
              gradiente + dados + computação
      \end{itemize}
    \end{column}
    \begin{column}{0.34\textwidth}
      \fotoslot{nobel-2024.jpg}{3.2cm}{2.35cm}
        {J.\ Hopfield e G.\ Hinton — coletiva do Nobel de Física 2024}
        {Foto: A.\ Petron, Wikimedia Commons, CC BY-SA 4.0}
    \end{column}
  \end{columns}

  \vspace{-0.55em}
  \begin{alertabox}
    A receita é simples de enunciar e difícil de dominar. Este curso existe
    para você entender \emph{cada peça} — a ponto de implementá-la do zero.
  \end{alertabox}
\end{frame}

% ------------------------------------------------------------
\begin{frame}{Por que só agora?}
  \framesubtitle{As ideias centrais são dos anos 1950--1980; o que mudou foi a escala}

  \begin{columns}[T]
    \begin{column}{0.31\textwidth}
      \begin{tcolorbox}[iccaixa=icNeural, title={1. Dados}]
        \small Conjuntos maiores; coleta e armazenamento baratos.
      \end{tcolorbox}
    \end{column}
    \begin{column}{0.31\textwidth}
      \begin{tcolorbox}[iccaixa=icAzulClaro, title={2. Hardware}]
        \small GPUs: a operação central — multiplicação de matrizes — é
        massivamente paralelizável.
      \end{tcolorbox}
    \end{column}
    \begin{column}{0.31\textwidth}
      \begin{tcolorbox}[iccaixa=icMagenta, title={3. Software}]
        \small Técnicas melhores, modelos novos, bibliotecas maduras.
      \end{tcolorbox}
    \end{column}
  \end{columns}

  \medskip
  \begin{center}\small
    {\color{icAzul}\textbf{1952}} gradiente descendente estocástico
    \;\textbullet\;
    {\color{icAzul}\textbf{1958}} perceptron com pesos aprendidos
    \;\textbullet\;
    {\color{icAzul}\textbf{1986}} retropropagação
    \;\textbullet\;
    {\color{icAzul}\textbf{1995}} redes convolucionais profundas
  \end{center}

  \begin{alertabox}
    As ideias centrais são dos anos 1950--1980. O que mudou: a escala — e
    escala muda o que é possível, não apenas o quão rápido.
  \end{alertabox}
\end{frame}

% ------------------------------------------------------------
\begin{frame}{O que uma rede profunda fez pela biologia}
  \framesubtitle{AlphaFold: o problema de 50 anos da estrutura de proteínas}

  \begin{columns}[T]
    \begin{column}{0.60\textwidth}
      \begin{center}
      \scalebox{0.58}{%
      \begin{tikzpicture}[node distance=0.55cm and 0.85cm]
        \node[bloco fuzzy, font=\ttfamily\small] (seq) {MKVLLA\ldots GRT};
        \node[below=1pt of seq, font=\scriptsize, text=icCinza] {sequência de aminoácidos};
        \node[bloco, right=of seq, align=center] (rede) {rede profunda\\ (atenção + geometria)};
        \node[right=1.1cm of rede] (prot) {};
        \draw[icAzulClaro, very thick,
              decorate, decoration={coil, aspect=0.62, segment length=3.4pt, amplitude=3.2pt}]
          ([xshift=0.15cm, yshift=-0.32cm]prot.west) .. controls +(0.55,0.75) and +(-0.5,0.55) ..
          ++(1.55,0.12);
        \draw[icMagenta, very thick]
          ([xshift=0.32cm, yshift=-0.42cm]prot.west) .. controls +(0.9,-0.42) and +(-0.35,-0.62) ..
          ++(1.5,0.28);
        \node[below=0.5cm of prot.east, xshift=-0.15cm, font=\scriptsize, text=icCinza]
          {estrutura 3D prevista};
        \draw[fluxo] (seq) -- (rede);
        \draw[fluxo] (rede) -- ([xshift=0.2cm]prot.west |- rede.east);
      \end{tikzpicture}}
      \end{center}

      \begin{itemize}\setlength{\itemsep}{0.25ex}
        \item Em 2020 (CASP14), atingiu precisão comparável à experimental —
              problema aberto desde os anos 1970
        \item Banco público com centenas de milhões de proteínas; Nobel de
              \emph{Química} 2024 para Hassabis e Jumper
      \end{itemize}
    \end{column}
    \begin{column}{0.36\textwidth}
      \fotoslot{alphafold-tmem253.jpg}{4.3cm}{2.2cm}
        {Predição do AlphaFold, colorida pela confiança}
        {Wikimedia Commons, CC0}
    \end{column}
  \end{columns}

  \vspace{-0.75em}
  \begin{ideiabox}
    Por dentro: modelo diferenciável de ponta a ponta, treinado por
    gradiente — feito de \emph{atenção}, onde o curso termina.
  \end{ideiabox}
\end{frame}

% ------------------------------------------------------------
\begin{frame}{Como o semestre funciona}
  \begin{columns}[T]
    \begin{column}{0.48\textwidth}
      \begin{tcolorbox}[iccaixa=icCiano, title={Em aula}]
        \begin{itemize}\setlength{\itemsep}{0.2ex}
          \item Teoria \emph{com} código: toda ideia central em PyTorch
          \item Perguntas para discutir — pense antes do clique da resposta
          \item Simuladores interativos para levar para casa
        \end{itemize}
      \end{tcolorbox}
    \end{column}
    \begin{column}{0.48\textwidth}
      \begin{tcolorbox}[iccaixa=icMagenta, title={Fora de aula}]
        \begin{itemize}\setlength{\itemsep}{0.2ex}
          \item Leituras: Goodfellow et al.\ (cap.\ 1, 6), \emph{d2l.ai},
                Nielsen
          \item Projeto: reprodução de um artigo, em grupos, com código
                rodando de verdade
        \end{itemize}
      \end{tcolorbox}
    \end{column}
  \end{columns}

  \medskip
  \begin{itemize}
    \item Material-base: CMU 11-785 (S26), MIT 6.S191, Stanford CS231n —
          adaptados e traduzidos para a nossa ementa
  \end{itemize}
\end{frame}

% ============================================================
\section{Conexionismo: a aposta que demorou a dar certo}

% ------------------------------------------------------------
\begin{frame}{Duas visões de como fazer uma máquina pensar}
  \begin{columns}[T]
    \begin{column}{0.48\textwidth}
      \begin{tcolorbox}[iccaixa=icCinza, title={IA simbólica}]
        \begin{itemize}\setlength{\itemsep}{0.2ex}
          \item Inteligência = manipular \emph{símbolos} com regras
          \item Conhecimento escrito à mão por especialistas
          \item Frágil fora do previsto
        \end{itemize}
      \end{tcolorbox}
    \end{column}
    \begin{column}{0.48\textwidth}
      \begin{tcolorbox}[iccaixa=icNeural, title={Conexionismo}]
        \begin{itemize}\setlength{\itemsep}{0.2ex}
          \item Inteligência \emph{emerge} de unidades simples conectadas
          \item Conhecimento fica nos \alert{pesos} — \emph{aprendido}
                de dados
          \item Inspiração: o cérebro ($\sim 10^{11}$ neurônios,
                $\sim 10^{14}$ sinapses)
        \end{itemize}
      \end{tcolorbox}
    \end{column}
  \end{columns}

  \vspace{-0.25em}
  \begin{ideiabox}
    A pergunta de hoje: unidades tão simples — somar e comparar com um limiar
    — \emph{conseguem} computar algo interessante? Sim — o argumento é
    concreto.
  \end{ideiabox}
\end{frame}

% ------------------------------------------------------------
\begin{frame}{Linha do tempo — 80 anos em um slide}
  \begin{itemize}\setlength{\itemsep}{0.2ex}
    \item<1-> {\color{icAzul}\textbf{1943}} McCulloch \& Pitts: o neurônio como
          \emph{unidade de limiar} — redes de neurônios computam lógica
    \item<1-> {\color{icAzul}\textbf{1949}} Hebb: sinapses que se reforçam com o
          uso — aprendizado local, sem professor
    \item<2-> {\color{icAzul}\textbf{1958}} Rosenblatt: o \alert{Perceptron} —
          pela primeira vez, os pesos são \emph{aprendidos} de exemplos
    \item<2-> {\color{icCarmim}\textbf{1969}} Minsky \& Papert: limites do
          perceptron simples (XOR) $\Rightarrow$ primeiro ``inverno''
    \item<3-> {\color{icAzul}\textbf{1986}} Rumelhart, Hinton \& Williams:
          retropropagação populariza o treino de redes \emph{multicamadas}
    \item<3-> {\color{icAzul}\textbf{2012}} AlexNet: GPUs + dados —
          deep learning vence em visão e nunca mais sai
    \item<4-> {\color{icMagenta}\textbf{2017--hoje}} o Transformer
          (\emph{atenção}) unifica linguagem, visão e áudio — é onde este
          curso quer chegar
  \end{itemize}
\end{frame}

% ------------------------------------------------------------
\begin{frame}{Os rostos desta história}
  \framesubtitle{Pitts, Rosenblatt, Widrow, Minsky, Papert, Rumelhart — os marcos de 1943 a 1986}

  \begin{columns}[T]
    \begin{column}{0.32\textwidth}
      \fotoslot{pitts-1954.jpg}{2.05cm}{1.3cm}
        {W.\ Pitts, 1954}
        {F.\ Bello, domínio público}
      \fotoslot{rosenblatt.jpg}{2.05cm}{1.3cm}
        {F.\ Rosenblatt}
        {Wikimedia Commons, CC BY-SA 4.0}
    \end{column}
    \begin{column}{0.32\textwidth}
      \fotoslot{widrow-adaline.jpg}{2.05cm}{1.3cm}
        {B.\ Widrow -- ADALINE}
        {Stanford Today, 1963, domínio público}
      \fotoslot{minsky.jpg}{2.05cm}{1.3cm}
        {M.\ Minsky}
        {S.\ Woodworth, CC BY 3.0}
    \end{column}
    \begin{column}{0.32\textwidth}
      \fotoslot{papert.jpg}{2.05cm}{1.3cm}
        {S.\ Papert}
        {Matematicamente.it, CC BY-SA 3.0}
      \fotoslot{rumelhart.jpg}{2.05cm}{1.3cm}
        {D.\ Rumelhart}
        {CC BY-SA 4.0}
    \end{column}
  \end{columns}

  \vspace{-0.35em}
  \begin{ideiabox}
    Dois marcos esquecidos: o ADALINE (1960) e a regra delta de Widrow--Hoff
    — a avó da descida de gradiente da aula 2.
  \end{ideiabox}
\end{frame}

% ============================================================
\section{O perceptron por dentro}

% ------------------------------------------------------------
\begin{frame}{O desenho que começou tudo}
  \framesubtitle{Cajal mostrou o neurônio como célula individual — McCulloch e Pitts o traduziram em matemática}

  \begin{columns}[T]
    \begin{column}{0.34\textwidth}
      \fotoslot{cajal-purkinje.jpg}{4.1cm}{2.8cm}
        {Purkinje (cerebelo humano) — desenho de Cajal}
        {Cajal Institute (CSIC), Madri — domínio público}
    \end{column}
    \begin{column}{0.62\textwidth}
      \begin{itemize}\setlength{\itemsep}{0.25ex}
        \item \textbf{1906:} Cajal, Nobel de Medicina pela
              \alert{doutrina do neurônio} — células individuais que se
              comunicam por \alert{sinapses}
        \item \textbf{1943:} McCulloch \& Pitts: somar entradas ponderadas e
              comparar com um limiar
        \item Do desenho enciclopédico sobra um modelo de poucas linhas — e
              é ele que sustenta o deep learning
      \end{itemize}
      \vspace{-0.7em}
      \begin{ideiabox}
        A história da área: sucessão de \emph{simplificações felizes} —
        descartam detalhe biológico, ganham poder de cálculo.
      \end{ideiabox}
    \end{column}
  \end{columns}
\end{frame}

% ------------------------------------------------------------
\begin{frame}{O original biológico}
  \framesubtitle{O neurônio de verdade — e o mapa que McCulloch e Pitts fizeram dele}

  \begin{columns}[T]
    \begin{column}{0.52\textwidth}
      \begin{center}
      \scalebox{0.74}{%
      \begin{tikzpicture}[x=0.82cm, y=0.82cm]
        % dendritos
        \foreach \a/\b in {150/1.55, 175/1.7, 205/1.6, 125/1.35}{
          \draw[icCiano, thick] (0,0) -- (\a:\b)
            \foreach \d in {-22,22}{ -- +(\a+\d:0.45) -- +(\a:0)} ;}
        % soma
        \fill[icAzul!14] (0,0) circle (0.42);
        \draw[icAzul, thick] (0,0) circle (0.42);
        \node[font=\scriptsize, text=icAzul] at (0,0) {soma};
        % axonio
        \draw[icCinza!85, very thick] (0.42,0) -- (2.6,0);
        \foreach \d in {28,-28,0}{
          \draw[icMagenta, thick] (2.6,0) -- +(\d:0.55);}
        % rotulos
        \node[font=\scriptsize, text=icCiano, align=center] at (-1.75,1.35)
          {dendritos\\ (entradas)};
        \node[font=\scriptsize, text=icCinza] at (1.5,0.3) {axônio};
        \node[font=\scriptsize, text=icMagenta, align=center] at (3.15,-0.85)
          {terminais\\ sinápticos (saída)};
      \end{tikzpicture}}
      \end{center}
      {\scriptsize Esquema; o neurônio real tem $10^{3}$--$10^{4}$ sinapses
      de entrada.\par}
    \end{column}
    \begin{column}{0.44\textwidth}
      \begin{itemize}\setlength{\itemsep}{0.35ex}
        \item Dendritos $\to$ \alert{entradas} $x_i$
        \item Eficácia da sinapse $\to$ \alert{peso} $w_i$
        \item Soma integra $\to$ $\vw^{\!\top}\vx$
        \item Limiar de disparo $\to$ \alert{ativação} $\ativ$
      \end{itemize}
      \begin{ideiabox}
        A abstração de 1943 mantém \emph{só} somar e comparar com limiar —
        e ainda assim basta para computar qualquer coisa.
      \end{ideiabox}
    \end{column}
  \end{columns}
\end{frame}

% ------------------------------------------------------------
\begin{frame}{O modelo}
  \begin{columns}[T]
    \begin{column}{0.52\textwidth}
      \begin{defbox}{Neurônio de limiar (perceptron)}
        Campo local induzido \(\campo = \vw^{\!\top}\vx + b\); saída
        \(\hat{y} = \ativ(\campo)\), com \(\ativ\) o degrau:
        \(\ativ(\campo) = 1\) se \(\campo \ge 0\), senão \(0\).
      \end{defbox}
      \begin{itemize}\setlength{\itemsep}{0.2ex}
        \item $\vw$: quanto cada entrada \emph{importa}
        \item $b$: quão fácil é disparar
        \item $\ativ$: a decisão (não linearidade)
      \end{itemize}
    \end{column}
    \begin{column}{0.44\textwidth}
      \begin{center}
        \scalebox{0.72}{\perceptron}
      \end{center}
    \end{column}
  \end{columns}

  \begin{ideiabox}
    Geometria: \(\vw^{\!\top}\vx + b = 0\) é um \emph{hiperplano}; $\vw$ é a
    normal, $b$ desloca. O perceptron é um separador linear — e essa frase
    inteira já contém a sua limitação.
  \end{ideiabox}
\end{frame}

% ------------------------------------------------------------
\begin{frame}{Do vetor ao lote: a forma matricial}
  \framesubtitle{A mesma conta, de uma entrada por vez a um lote inteiro}

  \[
    \hat{y} = \ativ(w_0 + \vx^{\!\top}\vw)
    \qquad\Longrightarrow\qquad
    Z = XW + b, \qquad A = \ativ(Z)
  \]

  \begin{itemize}\setlength{\itemsep}{0.35ex}
    \item O viés $w_0$ desloca a fronteira de decisão — sem ele, ela passa
          pela origem
    \item Lote de $B$ exemplos: $X$ é $B{\times}m$, $W$ é $m{\times}1$ e a
          saída, $B{\times}1$ — o lote atravessa a conta intocado
    \item Perceptron = produto interno + função escalar — a operação que as
          GPUs fazem bem
  \end{itemize}

  \begin{ideiabox}
    Rede profunda: pilha de multiplicações de matrizes intercaladas com não
    linearidades. Erros de dimensão: bug nº 1 — colunas da entrada casam com
    linhas dos pesos, sempre.
  \end{ideiabox}
\end{frame}

% ------------------------------------------------------------
\begin{frame}{A geometria da decisão}
  \framesubtitle{Quatro pontos bastam para ver o que o perceptron pode — e o que não pode}

  \begin{columns}[T]
    \begin{column}{0.31\textwidth}
      \centering {\color{icVerde}\bfseries\small E: separável}\\[2pt]
      \begin{tikzpicture}[scale=1.05]
        \eixosbool
        \pzero{0}{0}\pzero{0}{1.3}\pzero{1.3}{0}
        \pum{1.3}{1.3}
        \draw[icVerde, very thick, dashed] (0.55,1.75) -- (1.75,0.55);
      \end{tikzpicture}
    \end{column}
    \begin{column}{0.31\textwidth}
      \centering {\color{icVerde}\bfseries\small OU: separável}\\[2pt]
      \begin{tikzpicture}[scale=1.05]
        \eixosbool
        \pzero{0}{0}
        \pum{0}{1.3}\pum{1.3}{0}\pum{1.3}{1.3}
        \draw[icVerde, very thick, dashed] (-0.3,0.95) -- (0.95,-0.3);
      \end{tikzpicture}
    \end{column}
    \begin{column}{0.31\textwidth}
      \centering {\color{icCarmim}\bfseries\small XOR: impossível}\\[2pt]
      \begin{tikzpicture}[scale=1.05]
        \eixosbool
        \pzero{0}{0}\pzero{1.3}{1.3}
        \pum{0}{1.3}\pum{1.3}{0}
        \draw[icCinza!55, thick, dashed] (-0.3,0.95) -- (1.75,0.6);
        \draw[icCinza!55, thick, dashed] (0.5,-0.35) -- (0.85,1.8);
        \node[text=icCarmim, font=\Large\bfseries] at (1.62,1.62) {?};
      \end{tikzpicture}
    \end{column}
  \end{columns}

  \medskip
  \begin{itemize}\setlength{\itemsep}{0.2ex}
    \item {\color{icMagenta}$\bullet$}~classe 1 \quad
          {\color{icCiano}$\circ$}~classe 0 \quad — uma reta
          $\vw^{\!\top}\vx + b = 0$ decide tudo
    \item No XOR, qualquer reta deixa um ponto do lado errado: não é o
          algoritmo de treino, é a \alert{família de funções}
  \end{itemize}
\end{frame}

% ------------------------------------------------------------
\begin{frame}{Perceptrons são portas lógicas}
  \framesubtitle{O argumento de McCulloch--Pitts, em três linhas de pesos}

  \begin{columns}[T]
    \begin{column}{0.55\textwidth}
      Com entradas booleanas $x_i \in \{0,1\}$ e degrau em $\campo \ge 0$:
      \begin{itemize}\setlength{\itemsep}{0.35ex}
        \item \textbf{E:}\; $\vw = (1,1)$, $b = -1{,}5$
              \;$\Rightarrow$\; dispara só com $x_1 = x_2 = 1$
        \item \textbf{OU:}\; $\vw = (1,1)$, $b = -0{,}5$
        \item \textbf{NÃO:}\; $w = -1$, $b = 0{,}5$
      \end{itemize}

      \smallskip
      \begin{teobox}{Universalidade booleana}
        \{E, OU, NÃO\} é um conjunto universal. Logo, \emph{redes} de
        perceptrons computam \textbf{qualquer função booleana} — qualquer
        circuito digital, qualquer tabela-verdade.
      \end{teobox}
    \end{column}
    \begin{column}{0.41\textwidth}
      \begin{center}
        \scalebox{0.95}{%
        \begin{tikzpicture}
          \node[neuronio entrada] (a) {$x_1$};
          \node[neuronio entrada, below=0.85cm of a] (b) {$x_2$};
          \node[neuronio, right=1.8cm of a, yshift=-0.62cm] (e) {E};
          \draw[sinapse] (a) -- node[peso] {$1$} (e);
          \draw[sinapse] (b) -- node[peso] {$1$} (e);
          \node[neuronio saida, right=1.5cm of e] (s) {$\hat{y}$};
          \draw[sinapse destac] (e) -- (s);
        \end{tikzpicture}}
      \end{center}
      \small O neurônio é um \emph{votante com limiar}; a lógica é o caso
      particular em que os votos são inteiros.
    \end{column}
  \end{columns}
\end{frame}

% ------------------------------------------------------------
\begin{frame}{Rosenblatt: os pesos podem ser aprendidos}
  \begin{algbox}{Regra do perceptron (1958)}
    Para cada exemplo $(\vx, y)$ com $y \in \{0,1\}$: calcule
    $\hat{y} = \ativ(\vw^{\!\top}\vx + b)$ e, se errou, corrija na direção
    do erro:
    \[
      \vw \leftarrow \vw + \termo{\taxa}{taxa de aprendizado}
                     \,(\,\termo{y - \hat{y}}{erro: $-1$, $0$ ou $+1$}\,)\,\vx
      \qquad
      b \leftarrow b + \taxa\,(y - \hat{y})
    \]
  \end{algbox}

  \begin{teobox}{Convergência do perceptron (Novikoff, 1962)}
    Se os dados são \alert{linearmente separáveis}, a regra converge em um
    número finito de atualizações. Se não são — ela oscila para sempre.
  \end{teobox}

  \begin{itemize}
    \item A forma \(\vw \leftarrow \vw + \Delta\vw\): \emph{ancestral
          direto} da descida de gradiente do semestre todo
  \end{itemize}
\end{frame}

% ------------------------------------------------------------
\begin{frame}{1958, em hardware: o Mark I Perceptron}
  \framesubtitle{Antes de ser uma linha de PyTorch, o peso era um botão girado por um motor}

  \begin{columns}[T]
    \begin{column}{0.46\textwidth}
      \fotoslot{mark1.jpg}{4.6cm}{2.75cm}
        {O Mark I Perceptron, do Cornell Aeronautical Laboratory}
        {Foto: National Museum of the U.S. Navy — domínio público}
    \end{column}
    \begin{column}{0.5\textwidth}
      \begin{itemize}\setlength{\itemsep}{0.35ex}
        \item ``Retina'' de $20{\times}20$ fotocélulas: $\vx$ com 400
              dimensões
        \item Cada peso $w_i$: um \alert{potenciômetro}; a regra girava
              os botões com motores elétricos
        \item Em 1958, a imprensa prometeu máquinas que andariam e
              falariam
      \end{itemize}
      \begin{ideiabox}
        \emph{Hype} seguido de decepção não é invenção da nossa década —
        o ciclo 1958--1969 calibra a leitura das manchetes de hoje.
      \end{ideiabox}
    \end{column}
  \end{columns}
\end{frame}

% ------------------------------------------------------------
\begin{frame}[fragile]{A regra do perceptron em PyTorch}
  \framesubtitle{Dez linhas — aprendendo a porta E}

\begin{lstlisting}
import torch

X = torch.tensor([[0.,0.],[0.,1.],[1.,0.],[1.,1.]])
y = torch.tensor([0., 0., 0., 1.])            # porta E

w, b = torch.zeros(2), torch.zeros(1)
for epoca in range(10):
    for xi, yi in zip(X, y):
        y_hat = (w @ xi + b >= 0).float()     # degrau
        erro  = yi - y_hat                    # -1, 0 ou +1
        w += erro * xi                        # regra do perceptron
        b += erro
\end{lstlisting}

  \begin{itemize}\setlength{\itemsep}{0.2ex}
    \item Sem \texttt{backward()}, sem gradiente: a regra é \emph{ad hoc} —
          degrau + erro discreto conspiram a favor
    \item Troque a porta E por XOR na linha 4 e rode: \alert{nunca converge}.
          Por quê?
  \end{itemize}
\end{frame}

% ------------------------------------------------------------
\begin{frame}{Para discutir em aula}
  \begin{quizbox}
    O perceptron simples aprende E e OU sem dificuldade. Então, com paciência e
    uma taxa de aprendizado bem escolhida, ele também acaba aprendendo XOR
    ($x_1 \oplus x_2$).

    \smallskip
    \opcoes{Verdadeiro \;/\; Falso}
  \end{quizbox}

  \pause

  \begin{respbox}
    \textbf{Falso — e não é questão de paciência.} Vimos a geometria:
    \alert{nenhuma reta} separa $\{(0,0),(1,1)\}$ de $\{(0,1),(1,0)\}$, e o
    teorema de convergência exigia exatamente separabilidade. Foi este
    exemplo, em Minsky \& Papert (1969), que congelou a área por uma década.
    A saída não é um perceptron melhor: é \textbf{empilhar} perceptrons.
  \end{respbox}
\end{frame}

% ============================================================
\section{Profundidade: de portas a aproximadores universais}

% ------------------------------------------------------------
\begin{frame}{Uma camada oculta resolve o XOR}
  \framesubtitle{XOR $=$ (OU) E NÃO (E) — a rede escreve essa fórmula em pesos}

  \begin{columns}[T]
    \begin{column}{0.5\textwidth}
      \begin{center}
        \scalebox{0.98}{%
        \begin{tikzpicture}
          \node[neuronio entrada] (x1) {$x_1$};
          \node[neuronio entrada, below=1.05cm of x1] (x2) {$x_2$};
          \node[neuronio, right=2cm of x1] (h1) {OU};
          \node[neuronio, right=2cm of x2] (h2) {E};
          \node[neuronio saida, right=5.3cm of x1, yshift=-0.8cm] (y) {$\hat{y}$};
          \draw[sinapse] (x1) -- node[peso] {$1$} (h1);
          \draw[sinapse] (x2) -- node[peso] {$1$} (h1);
          \draw[sinapse] (x1) -- node[peso] {$1$} (h2);
          \draw[sinapse] (x2) -- node[peso] {$1$} (h2);
          \draw[sinapse destac] (h1) -- node[peso] {$1$} (y);
          \draw[sinapse destac] (h2) -- node[peso] {$-2$} (y);
        \end{tikzpicture}}
      \end{center}
      \small Limiares: $h_1$ dispara com $\campo \ge 0{,}5$; $h_2$ com
      $\campo \ge 1{,}5$; $\hat y$ com $\campo \ge 0{,}5$. Verifique nas
      quatro entradas.
    \end{column}
    \begin{column}{0.46\textwidth}
      \begin{itemize}\setlength{\itemsep}{0.35ex}
        \item $h_1$ e $h_2$ são \emph{features} construídas: ``algum ligado''
              e ``ambos ligados''
        \item Na saída, uma única reta sobre $(h_1, h_2)$ resolve — o
              próximo slide mostra isso acontecendo
      \end{itemize}
      \begin{ideiabox}
        A tese do deep learning inteiro: camadas não adicionam só capacidade —
        adicionam \textbf{representações}.
      \end{ideiabox}
    \end{column}
  \end{columns}
\end{frame}

% ------------------------------------------------------------
\begin{frame}{O que a camada oculta faz com o espaço}
  \framesubtitle{Os mesmos quatro pontos, antes e depois de $(h_1, h_2)$}

  \begin{columns}[T]
    \begin{column}{0.34\textwidth}
      \centering {\bfseries\small Espaço de entrada $(x_1, x_2)$}\\[2pt]
      \begin{tikzpicture}[scale=1.12]
        \eixosbool
        \pzero{0}{0}\pzero{1.3}{1.3}
        \pum{0}{1.3}\pum{1.3}{0}
        \node[text=icCarmim, font=\bfseries\small] at (1.55,1.7) {?};
      \end{tikzpicture}
    \end{column}
    \begin{column}{0.16\textwidth}
      \vspace{2.1cm}
      \centering
      {\Large $\xrightarrow{\;h\;}$}\\[2pt]
      {\scriptsize\color{icCinza} camada oculta}
    \end{column}
    \begin{column}{0.34\textwidth}
      \centering {\bfseries\small Espaço oculto $(h_1, h_2)$}\\[2pt]
      \begin{tikzpicture}[scale=1.12]
        \draw[-{Stealth}, icCinza!80] (-0.45,0) -- (1.9,0) node[below, font=\scriptsize] {$h_1$};
        \draw[-{Stealth}, icCinza!80] (0,-0.45) -- (0,1.9) node[left, font=\scriptsize] {$h_2$};
        \node[font=\tiny, text=icCinza] at (1.3,-0.25) {$1$};
        \node[font=\tiny, text=icCinza] at (-0.25,1.3) {$1$};
        \pzero{0}{0}\pzero{1.3}{1.3}
        \pum{1.3}{0}
        \node[font=\tiny, text=icCinza] at (1.3,-0.55) {$(0,1)$ e $(1,0)$ caem juntos};
        \draw[icVerde, very thick, dashed] (0.55,-0.35) -- (1.85,0.95);
      \end{tikzpicture}
    \end{column}
  \end{columns}

  \medskip
  \begin{itemize}\setlength{\itemsep}{0.2ex}
    \item A camada oculta \alert{dobra o espaço}: pontos que precisavam da
          mesma resposta caem no \emph{mesmo lugar}
    \item Depois, o problema volta a ser linear — e o perceptron de
          saída dá conta
  \end{itemize}
\end{frame}

% ------------------------------------------------------------
\begin{frame}{Funções de ativação comuns}
  \framesubtitle{Três suavizações do degrau — cada uma com sua derivada}

  \begin{columns}[T]
    \begin{column}{0.28\textwidth}
      \begin{tcolorbox}[iccaixa=icCiano, title={Sigmoide}]
        \small\[ \sig(z) = \frac{1}{1 + e^{-z}} \]
        \footnotesize\[ \sig'(z) = \sig(z)\bigl(1 - \sig(z)\bigr) \]
      \end{tcolorbox}
    \end{column}
    \begin{column}{0.28\textwidth}
      \begin{tcolorbox}[iccaixa=icAzulClaro, title={Tangente hiperbólica}]
        \small\[ \tanh(z) = \frac{e^{z} - e^{-z}}{e^{z} + e^{-z}} \]
        \footnotesize\[ \tanh'(z) = 1 - \tanh^{2}(z) \]
      \end{tcolorbox}
    \end{column}
    \begin{column}{0.28\textwidth}
      \begin{tcolorbox}[iccaixa=icMagenta, title={ReLU}]
        \small\[ \relu(z) = \max(0, z) \]
        \footnotesize $\relu'(z)$: $1$ se $z > 0$; $0$ caso contrário
      \end{tcolorbox}
    \end{column}
  \end{columns}

  \begin{itemize}\setlength{\itemsep}{0.3ex}
    \item Todas são \alert{não lineares} — e é exatamente esse o ponto
    \item Nas camadas ocultas de hoje, \alert{ReLU} é o padrão; na saída, a
          ativação depende da tarefa
  \end{itemize}
\end{frame}

% ------------------------------------------------------------
\begin{frame}{Para discutir em aula}
  \begin{quizbox}
    Uma rede de dez camadas com ativação identidade $g(z) = z$ tem o mesmo
    poder expressivo de uma única camada linear, não importa quantos
    neurônios cada camada tenha.

    \smallskip
    \opcoes{Verdadeiro \;/\; Falso}
  \end{quizbox}

  \pause

  \begin{respbox}
    \textbf{Verdadeiro.} Se $g$ é linear, a composição de camadas volta a ser
    uma função linear: \(y = W_2(W_1\vx + b_1) + b_2 = (W_2W_1)\vx + (W_2b_1 + b_2)\).
    Uma rede de 100 camadas teria o poder de uma só. A não linearidade é o que
    permite a cada camada acrescentar \alert{representação de verdade}.
  \end{respbox}
\end{frame}

% ------------------------------------------------------------
\begin{frame}{Propagação direta, passo a passo}
  \framesubtitle{Rede $2{\to}2{\to}1$, ReLU na camada oculta, sigmoide na saída — um exemplo numérico completo}

  \begin{columns}[T]
    \begin{column}{0.38\textwidth}
      \begin{center}
        \scalebox{0.66}{\mlp{2,2,1}}
      \end{center}
      \small Entrada \(\vx = (1, 2)\); pesos e viés:
      \[
        W^{(1)} = \begin{psmallmatrix} 0{,}5 & -1{,}0\\ 1{,}5 & 0{,}5 \end{psmallmatrix},\;
        b^{(1)} = (0, -1)
      \]
      \[
        W^{(2)} = \begin{psmallmatrix} 1{,}0\\ 2{,}0 \end{psmallmatrix},\;
        b^{(2)} = -2
      \]
    \end{column}
    \begin{column}{0.58\textwidth}
      \begin{center}\footnotesize
      \begin{tabular}{@{}cll@{}}
        \toprule
        \textbf{Passo} & \textbf{Operação} & \textbf{Resultado}\\
        \midrule
        1 & \(\vet{z}^{(1)} = \vx W^{(1)} + b^{(1)}\) & \((3{,}5,\; -1{,}0)\)\\
        2 & \(\vet{a}^{(1)} = \relu(\vet{z}^{(1)})\) & \((3{,}5,\; 0)\)\\
        3 & \(z^{(2)} = \vet{a}^{(1)} W^{(2)} + b^{(2)}\) & \(1{,}5\)\\
        4 & \(\hat{y} = \sig(z^{(2)})\) & \(0{,}8176\)\\
        \bottomrule
      \end{tabular}
      \end{center}
      \small Confira: \(3{,}5 = 1(0{,}5) + 2(1{,}5) + 0\) e
      \(\hat y = 1/(1 + e^{-1{,}5})\).
    \end{column}
  \end{columns}

  \begin{alertabox}
    Passo 2: a ReLU zerou a segunda unidade oculta — para esta entrada ela
    não contribui e, na retropropagação, não recebe gradiente. É assim que a
    ReLU produz \alert{ativações esparsas}.
  \end{alertabox}
\end{frame}

% ------------------------------------------------------------
\begin{frame}{Do booleano ao contínuo}
  \begin{columns}[T]
    \begin{column}{0.5\textwidth}
      \begin{center}
        \scalebox{0.52}{\mlp{2,4,1}}
      \end{center}
      \small Troque o degrau por uma $\ativ$ suave ($\sig$, tanh,
      $\relu$) e tudo vira diferenciável. Intuição: cada neurônio contribui
      um ``degrau''; degraus suficientes viram uma tabela de consulta suave.
    \end{column}
    \begin{column}{0.46\textwidth}
      \begin{teobox}{Aproximação universal (Cybenko 1989; Hornik 1991)}
        Uma rede com \emph{uma} camada oculta e ativação não polinomial
        aproxima qualquer função contínua em um compacto de $\R^n$.
      \end{teobox}
      \vspace{-0.25em}
      \begin{alertabox}
        O teorema diz que os pesos \emph{existem} — não quantos neurônios
        custam, nem como \emph{encontrá-los}. Representação, não aprendizado.
      \end{alertabox}
    \end{column}
  \end{columns}

\end{frame}

% ------------------------------------------------------------
\begin{frame}{A hierarquia que emerge sozinha}
  \framesubtitle{Ninguém programou ``detecte uma borda'' — as imagens abaixo são o que a rede aprendeu sozinha}

  \begin{columns}[T]
    \begin{column}{0.58\textwidth}
      \fotoslot{camadas-googlenet.png}{5.0cm}{3.2cm}
        {O que cada camada do GoogLeNet aprendeu sozinha: texturas, partes, padrões e objetos}
        {Olah, Mordvintsev \& Schubert — Feature Visualization, distill.pub, 2017, CC BY 4.0}
    \end{column}
    \begin{column}{0.38\textwidth}
      \begin{itemize}\setlength{\itemsep}{0.25ex}
        \item Cada imagem é um \alert{neurônio real} do GoogLeNet — ninguém
              programou ``detecte uma borda''
        \item A hierarquia \emph{emerge} do treinamento: cedo, texturas;
              fundo, partes e objetos
        \item O XOR em escala: lá, ``algum ligado'' e ``ambos ligados'';
              aqui, bordas, olhos e rostos
      \end{itemize}
    \end{column}
  \end{columns}
\end{frame}

% ------------------------------------------------------------
\begin{frame}{Por que fundo, e não largo?}
  \framesubtitle{O resultado mais importante da aula — e o mais fácil de perder}

  Paridade de $n$ bits, \(x_1 \oplus \dots \oplus x_n\):

  \begin{itemize}\setlength{\itemsep}{0.35ex}
    \item<1-> Com \textbf{uma} camada oculta, a construção por tabela-verdade
          usa \(\mdestaque{icCarmim}{2^{n-1}}\) neurônios; circuitos rasos de
          E/OU exigem tamanho exponencial (Håstad, 1986)
    \item<2-> Com \textbf{profundidade}: uma árvore de XORs de 2 entradas usa
          \(\mdestaque{icVerde}{\mathcal{O}(n)}\) neurônios em
          \(\mathcal{O}(\log n)\) camadas
  \end{itemize}

  \onslide<3->{%
  \vspace{-0.5em}
  \begin{ideiabox}
    Profundidade compra \textbf{reuso de computação intermediária}: cada camada
    calcula subexpressões que as seguintes combinam; a rede rasa trata caso a
    caso. É por isso que você escreve programas com funções, não uma tabela.
  \end{ideiabox}}

  \onslide<4->{%
  \vspace{-0.6em}
  \begin{itemize}
    \item \textbf{Fontes:} Raj, CMU 11-785, Lecs.\ 1--2; Håstad (1986);
          para redes contínuas, Telgarsky (2016) e Eldan \& Shamir (2016)
          — o ``deep'' é contagem de circuito, não moda
  \end{itemize}}
\end{frame}

% ------------------------------------------------------------
\begin{frame}{Para discutir em aula}
  \begin{quizbox}
    Pelo teorema da aproximação universal, uma camada oculta basta. Logo, redes
    profundas são um luxo de engenharia: qualquer coisa que uma rede profunda
    faz, uma rede rasa faz igualmente bem na prática.

    \smallskip
    \opcoes{Verdadeiro \;/\; Falso}
  \end{quizbox}

  \pause

  \begin{respbox}
    \textbf{Falso.} ``Existe uma rede rasa equivalente'' é verdade; ``do mesmo
    tamanho'' não é. Para famílias como a paridade, a versão rasa é
    \alert{exponencialmente} maior — impossível de armazenar, quanto mais de
    treinar. Universalidade é sobre \emph{o que é representável};
    profundidade é sobre \emph{a que custo}. O semestre vive nessa distância.
  \end{respbox}
\end{frame}

% ============================================================
\section{O que vem pela frente}

% ------------------------------------------------------------
\begin{frame}{A ativação de saída define a tarefa}
  \framesubtitle{Nas camadas ocultas, ReLU; na última camada, a escolha importa}

  \begin{center}\scriptsize
  \begin{tabular}{@{}lll@{}}
    \toprule
    \textbf{Tarefa} & \textbf{Ativação de saída} & \textbf{Interpretação}\\
    \midrule
    Regressão & nenhuma (linear) & valor real, sem restrição\\
    Classificação binária & sigmoide, 1 saída & $P(y = 1 \mid \vx)$\\
    Classificação multiclasse & softmax, $K$ saídas & distribuição sobre $K$ classes\\
    Multirrótulo & sigmoide, $K$ saídas & $K$ probabilidades independentes\\
    \bottomrule
  \end{tabular}
  \end{center}

  \vspace{-0.5em}
  \begin{defbox}{Softmax}
    \[
      \soft(\vet{z})_i = \frac{e^{z_i}}{\sum_{j=1}^{K} e^{z_j}}
      \qquad\Longrightarrow\qquad
      \sum_{i=1}^{K} \soft(\vet{z})_i = 1
    \]
  \end{defbox}

  \vspace{-0.5em}
  \begin{ideiabox}
    A sigmoide age em cada saída isoladamente; o softmax \emph{acopla} as
    saídas — subir uma classe baixa as outras.
  \end{ideiabox}
\end{frame}

% ------------------------------------------------------------
\begin{frame}[fragile]{Prévia: o XOR treinado por gradiente}
  \framesubtitle{O que a regra do perceptron não conseguiu, o gradiente resolve em 15 linhas}

\begin{lstlisting}
import torch, torch.nn as nn

X = torch.tensor([[0.,0.],[0.,1.],[1.,0.],[1.,1.]])
y = torch.tensor([[0.],[1.],[1.],[0.]])            # XOR

modelo = nn.Sequential(nn.Linear(2, 2), nn.Tanh(),
                       nn.Linear(2, 1), nn.Sigmoid())
otim  = torch.optim.SGD(modelo.parameters(), lr=0.5)
perda = nn.BCELoss()

for passo in range(5000):
    otim.zero_grad()
    L = perda(modelo(X), y)                # quao errado estou?
    L.backward()                           # retropropagacao
    otim.step()                            # w <- w - lr * grad
\end{lstlisting}

  \vspace{-0.4em}
  \begin{itemize}
    \item Repare: degrau $\to$ \texttt{Tanh}/\texttt{Sigmoid}. Sem
          suavidade não há derivada; sem derivada não há \texttt{backward()}
  \end{itemize}
\end{frame}

% ------------------------------------------------------------
\begin{frame}{A perda: medindo o erro}
  \framesubtitle{O que a \texttt{BCELoss} da prévia faz — e por que ela, e não o erro quadrático}

  \begin{defbox}{Perda empírica $J(\vw)$}
    \small \(J(\vw) = \frac{1}{n}\sum_{i=1}^{n}
    L\bigl(f(\vx^{(i)}; \vw),\, y^{(i)}\bigr)\) — função objetivo, custo ou
    risco empírico; treinar é procurar \(\argmin_{\vw} J(\vw)\).
  \end{defbox}

  \begin{columns}[T]
    \begin{column}{0.5\textwidth}
      \begin{tcolorbox}[iccaixa=icCiano, title={Entropia cruzada binária}]
        \footnotesize Classificação: a saída é uma probabilidade.
        \footnotesize\[ J = -\tfrac{1}{n}\sum_{i}
           \bigl[y_i\log \hat y_i + (1-y_i)\log(1-\hat y_i)\bigr] \]
      \end{tcolorbox}
    \end{column}
    \begin{column}{0.46\textwidth}
      \begin{tcolorbox}[iccaixa=icAzulClaro, title={Erro quadrático médio}]
        \footnotesize Regressão: a saída é um real contínuo.
        \footnotesize\[ J = \tfrac{1}{n}\sum_{i} \bigl(y_i - \hat y_i\bigr)^{2} \]
      \end{tcolorbox}
    \end{column}
  \end{columns}

  \begin{alertabox}
    EQM + sigmoide carrega $\sig'(z)$ no gradiente, que some quando o modelo
    está confiante e errado; a BCE cancela o $\sig'$ e sobra o erro
    $y - \hat y$.
  \end{alertabox}
\end{frame}

% ------------------------------------------------------------
\begin{frame}{As três linhas que valem o semestre}
  \[
    \texttt{L.backward()}
    \;\Longrightarrow\;
    \deriv{\erro}{\vw} \text{ para \emph{todos} os pesos, pela regra da cadeia}
  \]
  \[
    \texttt{otim.step()}
    \;\Longrightarrow\;
    \mdestaque{icMagenta}{\vw \leftarrow \vw - \taxa \deriv{\erro}{\vw}}
  \]

  \begin{itemize}\setlength{\itemsep}{0.35ex}
    \item Compare com a regra do perceptron: mesma forma
          \(\vw \leftarrow \vw + \Delta\vw\), mas agora \(\Delta\vw\) vem de um
          \alert{princípio} — minimizar uma perda diferenciável — e não de uma
          heurística
    \item Escala do XOR de 9 parâmetros aos modelos de linguagem com
          $10^{12}$: \emph{o algoritmo é o mesmo}
    \item Aulas 2--4: descida de gradiente, grafo computacional e a
          retropropagação implementada do zero (nosso mini-autograd)
  \end{itemize}
\end{frame}

% ------------------------------------------------------------
\begin{frame}{Mapa da disciplina}
  \begin{columns}[T]
    \begin{column}{0.4\textwidth}
      \begin{tcolorbox}[iccaixa=icNeural, title={Redes neurais e deep learning}]
        \small Perceptron $\to$ MLP $\to$ retropropagação e autograd $\to$
        otimização e regularização modernas $\to$ CNNs $\to$ sequências $\to$
        Hopfield moderno $\to$ atenção e Transformers
      \end{tcolorbox}
    \end{column}
    \begin{column}{0.28\textwidth}
      \begin{tcolorbox}[iccaixa=icEvol, title={Computação evolutiva}]
        \small Otimização \emph{sem} gradiente: AGs, estratégias evolutivas,
        CMA-ES. Quando a derivada não existe — e o custo de viver sem ela.
      \end{tcolorbox}
    \end{column}
    \begin{column}{0.28\textwidth}
      \begin{tcolorbox}[iccaixa=icFuzzy, title={Sistemas nebulosos}]
        \small Pertinência gradual $\to$ relaxações diferenciáveis: softmax,
        \emph{gating}, atenção como consulta ``fuzzy''.
      \end{tcolorbox}
    \end{column}
  \end{columns}

  \medskip
  \begin{itemize}\setlength{\itemsep}{0.25ex}
    \item O bloco {\color{icNeural}\textbf{neural}} carrega a maior parte
          da carga horária; os outros dois entram como \emph{contraponto}
    \item Ao final, você deve saber ler um artigo atual de deep learning —
          e reproduzi-lo
  \end{itemize}
\end{frame}

% ------------------------------------------------------------
\begin{frame}{Os paradigmas de aprendizado}
  \framesubtitle{A mesma rede e o mesmo gradiente — o que muda é de onde vem o sinal}

  \begin{center}\footnotesize
  \begin{tabular}{@{}lp{4.1cm}p{4.1cm}@{}}
    \toprule
    \textbf{Paradigma} & \textbf{Sinal de aprendizado} & \textbf{Exemplo de uso}\\
    \midrule
    Supervisionado & rótulo correto por exemplo & diagnóstico por imagem\\
    Não supervisionado & nenhum rótulo & segmentação de clientes\\
    Auto-supervisionado & rótulo derivado do próprio dado & pré-treino de modelos de linguagem\\
    Por reforço & recompensa escalar, possivelmente atrasada & jogos, robótica\\
    \bottomrule
  \end{tabular}
  \end{center}

  \begin{ideiabox}
    Nas quatro linhas, a rede é a mesma e o gradiente vem por retropropagação.
    Um modelo de linguagem moderno usa três em sequência: pré-treino
    auto-supervisionado, ajuste supervisionado e alinhamento por reforço
    (tema da disciplina GBC063).
  \end{ideiabox}
\end{frame}

% ------------------------------------------------------------
\begin{frame}{Simuladores interativos desta unidade}
  \framesubtitle{Para mexer em casa — cada um cabe numa página HTML, sem instalar nada}

  \begin{itemize}\setlength{\itemsep}{0.4ex}
    \item {\color{icCiano}\textbf{Playground do perceptron}} — arraste pontos
          no plano, veja a reta $\vw^{\!\top}\vx + b = 0$ se mover a cada
          atualização da regra; presets E, OU e XOR (para ver a oscilação
          eterna com os próprios olhos)
    \item {\color{icAzulClaro}\textbf{O espaço oculto do XOR}} — a
          transformação $(x_1,x_2) \to (h_1,h_2)$ animada, com controle dos
          pesos da camada oculta
    \item {\color{icAzul}\textbf{Aproximador universal 1D}} — some ``degraus''
          de neurônios sigmoide para caber numa função-alvo; contador de
          neurônios visível
    \item {\color{icMagenta}\textbf{Largo vs.\ fundo}} — paridade de $n$ bits:
          o mesmo problema resolvido raso e fundo, com o número de neurônios
          lado a lado
  \end{itemize}

  \begin{ideiabox}
    Regra da disciplina: se um conceito tem geometria, ele ganha um simulador.
  \end{ideiabox}
\end{frame}

% ------------------------------------------------------------
\begin{frame}{Para saber mais — sites e vídeos por assunto}
  \framesubtitle{Tudo gratuito e verificado; clique no link para abrir — ou copie do LMS}

  \begin{columns}[T]
    \begin{column}{0.5\textwidth}
      {\color{icCinza}\bfseries\scriptsize FUNDAMENTOS E HISTÓRIA}
      {\scriptsize
      \begin{itemize}\setlength{\itemsep}{0.3ex}
        \item \textbf{3Blue1Brown} — o que é uma rede neural? (vídeo):
              \href{https://youtu.be/aircAruvnKk}{\color{icAzulClaro}\texttt{youtu.be/aircAruvnKk}}
        \item \textbf{3Blue1Brown} — série completa (playlist):
              \href{https://www.youtube.com/playlist?list=PLZHQObOWTQDNU6R1_67000Dx_ZCJB-3pi}{\color{icAzulClaro}\texttt{youtube.com/playlist?list=PLZHQObO\ldots}}
        \item \textbf{MIT 6.S191} — introdução ao deep learning (curso):
              \href{https://introtodeeplearning.com}{\color{icAzulClaro}\texttt{introtodeeplearning.com}}
        \item \textbf{1958 em filme} — o Mark I aprende a ler (vídeo):
              \href{https://youtu.be/cNxadbrN_aI}{\color{icAzulClaro}\texttt{youtu.be/cNxadbrN\_aI}}
        \item \textbf{Nielsen} — cap.\ 1 e 4, provas visuais:
              \href{https://neuralnetworksanddeeplearning.com}{\color{icAzulClaro}\texttt{neuralnetworksanddeeplearning.com}}
        \item \textbf{Wikipedia} — Perceptron (pt):
              \href{https://pt.wikipedia.org/wiki/Perceptron}{\color{icAzulClaro}\texttt{pt.wikipedia.org/wiki/Perceptron}}
        \item \textbf{CMU 11-785} — palestras 1--2 (profundidade):
              \href{https://deeplearning.cs.cmu.edu}{\color{icAzulClaro}\texttt{deeplearning.cs.cmu.edu}}
      \end{itemize}}
    \end{column}
    \begin{column}{0.5\textwidth}
      {\color{icCinza}\bfseries\scriptsize DEEP LEARNING HOJE}
      {\scriptsize
      \begin{itemize}\setlength{\itemsep}{0.3ex}
        \item \textbf{Nobel 2024} — comunicado oficial (Física):
              \href{https://www.nobelprize.org/prizes/physics/2024/press-release/}{\color{icAzulClaro}\texttt{nobelprize.org/prizes/physics/2024}}
        \item \textbf{DeepMind} — AlphaFold, 50 anos resolvidos (blog):
              \href{https://deepmind.google/discover/blog/alphafold-a-solution-to-a-50-year-old-grand-challenge-in-biology/}{\color{icAzulClaro}\texttt{deepmind.google (blog AlphaFold)}}
        \item \textbf{PyTorch} — tutorial oficial:
              \href{https://pytorch.org/tutorials/beginner/basics/intro.html}{\color{icAzulClaro}\texttt{pytorch.org/tutorials}}
        \item \textbf{Karpathy} — Zero to Hero, do zero (curso):
              \href{https://karpathy.ai/zero-to-hero.html}{\color{icAzulClaro}\texttt{karpathy.ai/zero-to-hero.html}}
        \item \textbf{3Blue1Brown} — atenção e Transformers (vídeo):
              \href{https://youtu.be/eMlx5fFNoYc}{\color{icAzulClaro}\texttt{youtu.be/eMlx5fFNoYc}}
        \item \textbf{distill.pub} — o que cada camada aprende:
              \href{https://distill.pub/2017/feature-visualization/}{\color{icAzulClaro}\texttt{distill.pub/feature-visualization}}
      \end{itemize}}
    \end{column}
  \end{columns}
\end{frame}

% ------------------------------------------------------------
\begin{frame}{Para casa}
  \begin{exbox}{1 — Portas na mão}
    Encontre pesos e viés de um perceptron de degrau para NAND. Em seguida,
    usando só NANDs, desenhe uma rede de duas camadas que computa XOR e
    escreva os pesos de cada neurônio.
  \end{exbox}

  \begin{exbox}{2 — O fracasso, observado}
    Rode o código da regra do perceptron desta aula trocando a porta E por
    XOR. Registre $\vw$ e $b$ ao longo de 50 épocas e descreva o que acontece.
    Depois rode a versão com gradiente e compare.
  \end{exbox}

  \begin{itemize}\setlength{\itemsep}{0.25ex}
    \item \textbf{Leitura:} Goodfellow et al., cap.\ 1 e seções 6.1--6.4;
          Nielsen, cap.\ 1 (aproximação universal, visual)
    \item \textbf{Opcional:} 11-785 (S26), Lectures 1--2, para o argumento
          de profundidade vs.\ largura em detalhe
  \end{itemize}
\end{frame}

% ------------------------------------------------------------
\begin{frame}{Referências}
  \begin{columns}[T]
    \begin{column}{0.5\textwidth}
      {\color{icCinza}\bfseries\scriptsize FUNDADORES}
      {\scriptsize
      \begin{itemize}\setlength{\itemsep}{0.15ex}
        \item McCulloch, W.; Pitts, W.\ (1943). A logical calculus of the
              ideas immanent in nervous activity. \emph{Bull.\ Math.\
              Biophysics} 5.
        \item Rosenblatt, F.\ (1958). The perceptron: a probabilistic model
              [\ldots]. \emph{Psychological Review} 65(6).
        \item Novikoff, A.\ (1962). On convergence proofs on perceptrons.
              \emph{Symp.\ Math.\ Theory of Automata}.
        \item Minsky, M.; Papert, S.\ (1969). \emph{Perceptrons}. MIT Press.
        \item Rumelhart; Hinton; Williams (1986). Learning representations
              by back-propagating errors. \emph{Nature} 323.
        \item Cybenko (1989), \emph{MCSS} 2; Hornik (1991),
              \emph{Neural Networks} 4.
      \end{itemize}}
    \end{column}
    \begin{column}{0.5\textwidth}
      {\color{icCinza}\bfseries\scriptsize PROFUNDIDADE E HOJE}
      {\scriptsize
      \begin{itemize}\setlength{\itemsep}{0.15ex}
        \item Håstad (1986). Almost optimal lower bounds for small depth
              circuits. \emph{STOC}. (tb.\ Ajtai 1983; Furst, Saxe \& Sipser
              1984)
        \item Telgarsky (2016); Eldan \& Shamir (2016). \emph{COLT} —
              separações de profundidade para redes contínuas.
        \item Krizhevsky; Sutskever; Hinton (2012). ImageNet classification
              with deep CNNs. \emph{NeurIPS}.
        \item Jumper et al.\ (2021). Highly accurate protein structure
              prediction with AlphaFold. \emph{Nature} 596.
        \item Raj, B. \emph{11-785 Introduction to Deep Learning}, CMU
              (S26), Lecs.\ 1--2 — \texttt{deeplearning.cs.cmu.edu}
        \item Amini, A.; Amini, A. \emph{MIT 6.S191: Introduction to Deep
              Learning} (2026), \texttt{introtodeeplearning.com} — base
              complementar desta aula.
      \end{itemize}}
    \end{column}
  \end{columns}
  \smallskip
  \vspace{-0.3em}
  {\scriptsize \textbf{Livros:} Goodfellow, Bengio \& Courville (2016),
  \emph{Deep Learning}, MIT Press \textbullet\ Nielsen (2015),
  \texttt{neuralnetworksanddeeplearning.com}\par
  \textbf{Imagens:} Wikimedia Commons — A.\ Petron (CC BY-SA 4.0); S.\ Ramón
  y Cajal (domínio público); F.\ Bello (domínio público); F.\ Rosenblatt
  (CC BY-SA 4.0); Stanford Today (domínio público); S.\ Woodworth (CC BY 3.0);
  Matematicamente.it (CC BY-SA 3.0); D.\ Rumelhart (CC BY-SA 4.0); AlphaFold
  DB (CC0); Olah, Mordvintsev \& Schubert, \emph{Feature Visualization},
  distill.pub (CC BY 4.0); U.S.\ Navy (domínio público).\par}
\end{frame}

\end{document}
